% !TEX program = pdflatex
% =====================================================================
%  TwA Teaching Webinar — Emily's FERPA & AI segment (~20 min)
%  June 17, 2026.  Built from verified deep-research brief
%  (eb-lab/projects/twa-ra-bootcamp/2026-06-11_ferpa-ai-teaching-brief.md).
%  Preamble matches the shared UVM Econ design system.
% =====================================================================
\documentclass[aspectratio=169, 10pt]{beamer}

%% ============================================================
%% PACKAGES
%% ============================================================
\usepackage[T1]{fontenc}
\usepackage{lmodern}
\usepackage{graphicx}
\usepackage{booktabs}
\usepackage{array}
\usepackage{tabularx}
\usepackage{colortbl}
\usepackage{tikz}
\usetikzlibrary{positioning, calc, arrows.meta, shapes.geometric, fit, shadows}
\usepackage[most]{tcolorbox}
\usepackage{fontawesome5}
\usepackage{hyperref}
\usepackage{multicol}
\usepackage{xcolor}
\usepackage{textcomp}
\usepackage{pifont}

\graphicspath{{figures/}}

%% ============================================================
%% COLOR DEFINITIONS
%% ============================================================
\definecolor{UVMGreen}{HTML}{154734}
\definecolor{UVMGold}{HTML}{D4A84B}
\definecolor{SlateTeal}{HTML}{1B4D5C}
\definecolor{DeepCharcoal}{HTML}{2C3E50}
\definecolor{SoftCoral}{HTML}{C0534D}
\definecolor{CoolGray}{HTML}{8899A6}
\definecolor{PaleSage}{HTML}{EDF2F0}
\definecolor{LightGold}{HTML}{FDF6E3}
\definecolor{DarkSlate}{HTML}{0F2B36}
\definecolor{AccentBlue}{HTML}{3498DB}
\definecolor{GoGreen}{HTML}{2E7D52}

%% ============================================================
%% BEAMER THEME SETUP
%% ============================================================
\setbeamercolor{structure}{fg=SlateTeal}
\setbeamercolor{normal text}{fg=DeepCharcoal}
\setbeamercolor{frametitle}{fg=SlateTeal}
\setbeamercolor{title}{fg=SlateTeal}
\setbeamercolor{subtitle}{fg=UVMGold}
\setbeamercolor{author}{fg=DeepCharcoal}
\setbeamercolor{date}{fg=CoolGray}
\setbeamercolor{institute}{fg=CoolGray}
\setbeamercolor{footline}{fg=CoolGray}
\setbeamercolor{block title}{fg=white, bg=SlateTeal}
\setbeamercolor{block body}{bg=PaleSage}
\setbeamercolor{alerted text}{fg=SoftCoral}
\setbeamercolor{item}{fg=SlateTeal}
\setbeamercolor{subitem}{fg=CoolGray}

\setbeamerfont{title}{size=\Large, series=\bfseries}
\setbeamerfont{subtitle}{size=\normalsize}
\setbeamerfont{frametitle}{size=\large, series=\bfseries}
\setbeamerfont{footline}{size=\tiny}

\setbeamertemplate{navigation symbols}{}

% Custom frame title
\setbeamertemplate{frametitle}{%
  \vskip14pt%
  \nointerlineskip%
  {\usebeamerfont{frametitle}\usebeamercolor[fg]{frametitle}\insertframetitle\par}%
  \vskip2pt%
  {\color{UVMGold}\rule{\textwidth}{1.2pt}}%
  \vskip3pt%
}

% Custom footline
\setbeamertemplate{footline}{%
  \begin{beamercolorbox}[wd=\paperwidth, ht=16pt, right, rightskip=1.5cm]{footline}%
    \usebeamerfont{footline}\insertframenumber\,/\,\inserttotalframenumber\vskip3pt%
  \end{beamercolorbox}%
}

% Itemize styles
\setbeamertemplate{itemize item}{\color{SlateTeal}\raisebox{0.15ex}{\small$\blacktriangleright$}}
\setbeamertemplate{itemize subitem}{\color{CoolGray}\raisebox{0.1ex}{\footnotesize$\bullet$}}

% Compact tcolorbox defaults
\tcbset{
  keybox/.style={
    colback=PaleSage, colframe=SlateTeal, boxrule=0.8pt, arc=3pt,
    left=6pt, right=6pt, top=4pt, bottom=4pt,
    fonttitle=\bfseries\small\color{white}, coltitle=white, colbacktitle=SlateTeal,
  },
  accentbox/.style={
    colback=LightGold, colframe=UVMGold, boxrule=0.8pt, arc=3pt,
    left=6pt, right=6pt, top=4pt, bottom=4pt,
  },
  quotebox/.style={
    colback=PaleSage, colframe=SlateTeal!50, boxrule=0.5pt, arc=2pt,
    left=8pt, right=8pt, top=4pt, bottom=4pt,
    borderline west={3pt}{0pt}{UVMGold},
  },
  dobox/.style={
    colback=GoGreen!8, colframe=GoGreen, boxrule=0.8pt, arc=3pt,
    left=6pt, right=6pt, top=3pt, bottom=3pt,
    fonttitle=\bfseries\small\color{white}, coltitle=white, colbacktitle=GoGreen,
  },
  dontbox/.style={
    colback=SoftCoral!8, colframe=SoftCoral, boxrule=0.8pt, arc=3pt,
    left=6pt, right=6pt, top=3pt, bottom=3pt,
    fonttitle=\bfseries\small\color{white}, coltitle=white, colbacktitle=SoftCoral,
  },
  coralbox/.style={
    colback=SoftCoral!8, colframe=SoftCoral, boxrule=0.8pt, arc=3pt,
    left=6pt, right=6pt, top=3pt, bottom=3pt,
    fonttitle=\bfseries\small\color{white}, coltitle=white, colbacktitle=SoftCoral,
  },
}

% Citation helper (muted footnote-style cite)
\newcommand{\cmcite}[1]{{\scriptsize\color{CoolGray}\textit{#1}}}

%% ============================================================
%% METADATA
%% ============================================================
\title{FERPA \& AI Tools in Teaching}
\subtitle{What You Can and Can't Feed an AI}
\author{Emily Beam}
\institute{Thinking with Agents \;$\cdot$\; Teaching Webinar}
\date{June 17, 2026}

\begin{document}

%% ====== SLIDE 1: TITLE ======
{
\setbeamertemplate{footline}{}
\begin{frame}[plain]
\begin{tikzpicture}[remember picture, overlay]
  \fill[SlateTeal] (current page.north west) rectangle (current page.south east);
  \fill[UVMGold] ([yshift=0.3cm]current page.center -| current page.west) rectangle ([yshift=0.25cm]current page.center -| current page.east);
  \node[anchor=center, text=white, font=\fontsize{27}{33}\bfseries\selectfont]
    at ([yshift=2.2cm]current page.center) {FERPA \& AI Tools in Teaching};
  \node[anchor=center, text=UVMGold, font=\fontsize{14}{19}\selectfont]
    at ([yshift=1.25cm]current page.center) {What You Can --- and Can't --- Feed an AI};
  \node[anchor=center, text=white, font=\large\bfseries]
    at ([yshift=-0.6cm]current page.center) {Emily Beam};
  \node[anchor=center, text=PaleSage, font=\normalsize]
    at ([yshift=-1.3cm]current page.center) {University of Vermont \;$\cdot$\; Thinking with Agents};
  \node[anchor=center, text=CoolGray!70!white, font=\small]
    at ([yshift=-2.2cm]current page.center) {Teaching Webinar \quad$\cdot$\quad June 17, 2026};
  \node[text=white, opacity=0.04, font=\fontsize{100}{100}\selectfont, anchor=center, rotate=-15]
    at ([xshift=5cm, yshift=1.5cm]current page.center) {\faLock};
\end{tikzpicture}
\end{frame}
}

%% ====== SLIDE 2: THE 30-SECOND VERSION ======
\begin{frame}{The 30-Second Version}
\begin{tcolorbox}[accentbox, top=2pt, bottom=2pt]
\small
\textbf{FERPA applies the moment you put identifiable student work, grades, or education records into an AI tool.} There are exactly \textbf{two} clean, legal ways to use AI with student data.
\end{tcolorbox}
\vskip2pt
\begin{columns}[T]
\begin{column}{0.48\textwidth}
\begin{tcolorbox}[dobox, title={\faCheck\quad Path 1: Approved tool}, top=2pt, bottom=2pt]
\small Use an \textbf{institution-vetted} AI tool that qualifies as a ``school official'' (under contract / institutional control).
\end{tcolorbox}
\end{column}
\begin{column}{0.48\textwidth}
\begin{tcolorbox}[dobox, title={\faCheck\quad Path 2: De-identify}, top=2pt, bottom=2pt]
\small \textbf{Properly de-identify} the data first, so it's no longer protected information.
\end{tcolorbox}
\end{column}
\end{columns}
\vskip4pt
\begin{tcolorbox}[dontbox, title={\faBan\quad The one-line don't}, top=2pt, bottom=2pt]
\small Never paste identifiable student work or grades into your \textbf{personal/consumer} AI account (consumer ChatGPT, Claude, or Copilot). That's an unconsented disclosure of education records.
\end{tcolorbox}
\vskip1pt
\cmcite{Not legal advice --- accurate to 34 CFR Part 99 \& ED guidance as of June 2026; defer to your institution's policy.}
\end{frame}

%% ====== SECTION DIVIDER: FOUNDATION ======
{
\setbeamertemplate{footline}{}
\setbeamercolor{background canvas}{bg=SlateTeal}
\begin{frame}[plain, noframenumbering]
\centering\vfill
{\fontsize{22}{28}\bfseries\selectfont\color{white} Part 1}\\[6pt]
{\color{UVMGold}\rule{5cm}{0.8pt}}\\[6pt]
{\normalsize\color{UVMGold} The Legal Foundation}\\[4pt]
{\small\color{PaleSage} What FERPA Actually Requires}
\vfill
\end{frame}
}

%% ====== SLIDE 3: EDUCATION RECORDS ======
\begin{frame}{What's Protected: ``Education Records''}
\begin{columns}[T]
\begin{column}{0.54\textwidth}
\small Records that are:
\vskip3pt
\begin{enumerate}\setlength{\itemsep}{2pt}
  \item \textbf{directly related to a student}, AND
  \item \textbf{maintained by the institution} --- or by \textbf{a party acting for it}.
\end{enumerate}
\vskip5pt
Records an AI vendor holds \emph{on your institution's behalf} are \textbf{themselves} education records. The data stays under FERPA after you hand it over.
\vskip4pt
{\small Grades, transcripts, class lists, and student work are the canonical examples.}
\vskip4pt
\cmcite{34 CFR 99.3; ED ``Responsibilities of Third-Party Service Providers under FERPA'' (2015).}
\end{column}
\begin{column}{0.42\textwidth}
\begin{tcolorbox}[quotebox]
\small
\textbf{Useful nuance:}\\[3pt]
A \emph{student} pasting \emph{their own} work into a consumer tool does \textbf{not} make the vendor's records ``education records.''
\vskip4pt
The obligation runs to \textbf{you and the institution} --- not the student.
\end{tcolorbox}
\end{column}
\end{columns}
\end{frame}

%% ====== SLIDE 4: PII IS BROADER ======
\begin{frame}{What's Identifiable: PII Is Broader Than Names}
\begin{columns}[T]
\begin{column}{0.50\textwidth}
\begin{tcolorbox}[keybox, title={\faIdCard\quad FERPA's PII includes\ldots}]
\small
\begin{itemize}\setlength{\itemsep}{2pt}
  \item Direct identifiers (name, SSN, student ID)
  \item \textbf{Indirect} identifiers (DOB, place of birth)
  \item Interaction \textbf{metadata}
  \item Even \textbf{aggregate} data
\end{itemize}
\vskip3pt
\footnotesize \ldots anything that, alone or combined, lets ``a reasonable person in the school community identify the student with reasonable certainty.''
\end{tcolorbox}
\vskip3pt
\cmcite{34 CFR 99.3; ED Vendor FAQ.}
\end{column}
\begin{column}{0.46\textwidth}
\begin{tcolorbox}[coralbox, title={\faExclamationTriangle\quad The mosaic effect}]
\small
In a \textbf{small economics seminar}, an ``anonymized'' essay can still be re-identifiable from:
\vskip3pt
\begin{itemize}\setlength{\itemsep}{1pt}
  \item topic
  \item writing style
  \item other facts you already know
\end{itemize}
\vskip3pt
De-identification means addressing \textbf{combinations} of identifiers --- not just deleting names.
\end{tcolorbox}
\end{column}
\end{columns}
\end{frame}

%% ====== SLIDE 5: THE DEFAULT RULE — CONSENT ======
\begin{frame}{The Default Rule: Consent}
\begin{columns}[T]
\begin{column}{0.54\textwidth}
\small A school \textbf{may not disclose PII} from education records \textbf{without prior signed, dated written consent} --- \emph{unless} an enumerated exception applies.
\vskip5pt
``Disclosure'' explicitly covers electronic transmission ``by any means \ldots to any party.'' \textbf{Uploading to an external AI tool is a disclosure.}
\vskip4pt
\cmcite{34 CFR 99.30(a), 99.3; ED 2014 online-services guidance.}
\end{column}
\begin{column}{0.42\textwidth}
\begin{tcolorbox}[accentbox]
\small
\textbf{In higher ed, the consent is the \emph{student's} --- not the parent's.}
\vskip4pt
FERPA rights transfer to the student when they turn 18 \textbf{or enter a postsecondary institution at any age.}
\end{tcolorbox}
\vskip4pt
\begin{tcolorbox}[quotebox, top=3pt, bottom=3pt]
\scriptsize Being a ``disclosure'' $\neq$ automatically a violation --- it means an exception or de-identification must apply.
\end{tcolorbox}
\end{column}
\end{columns}
\end{frame}

%% ====== SLIDE 6: EXCEPTION A — SCHOOL OFFICIAL ======
\begin{frame}{Exception A: The ``School Official'' Pathway (for vendors)}
\vskip-2pt
\begin{columns}[T]
\begin{column}{0.55\textwidth}
\small A vendor may receive PII \textbf{without consent} only if it:
\vskip3pt
\begin{enumerate}\setlength{\itemsep}{2pt}
  \item performs a service the school would otherwise use \textbf{employees} for;
  \item is a designated school official with a \textbf{legitimate educational interest} (annual FERPA notice);
  \item is under the institution's \textbf{direct control}; and
  \item uses records \textbf{only for authorized purposes} --- \textbf{no re-disclosure}.
\end{enumerate}
\vskip4pt
\cmcite{34 CFR 99.31(a)(1)(i); 99.33(a); ED Vendor FAQ + 2014 guidance.}
\end{column}
\begin{column}{0.42\textwidth}
\begin{tcolorbox}[coralbox, title={\faTimesCircle\quad Bottom line}, top=3pt, bottom=3pt]
\small A \textbf{consumer/personal} AI account is \textbf{not} under institutional control. It \textbf{cannot} qualify.
\vskip3pt
Only an institutionally-contracted, vetted tool can.
\end{tcolorbox}
\vskip4pt
\begin{tcolorbox}[quotebox, top=3pt, bottom=3pt]
\scriptsize This same exception is \emph{why you} can lawfully see your own students' records --- instructors are ``school officials with legitimate educational interests.''
\end{tcolorbox}
\end{column}
\end{columns}
\end{frame}

%% ====== SLIDE 7: EXCEPTION B — DE-IDENTIFICATION ======
\begin{frame}{Exception B: De-Identification (no vendor contract needed)}
\begin{columns}[T]
\begin{column}{0.54\textwidth}
\small After removing all PII, an institution may use/release records without consent --- \emph{provided} it makes a \textbf{reasonable determination} that identity is not personally identifiable, accounting for single \emph{or multiple} releases and \textbf{other reasonably available information.}
\vskip4pt
Properly de-identified data is \textbf{outside FERPA entirely}, so it may be used in general AI tools.
\vskip4pt
\cmcite{34 CFR 99.31(b)(1); ED 2014 guidance.}
\end{column}
\begin{column}{0.42\textwidth}
\begin{tcolorbox}[coralbox, title={\faExclamationTriangle\quad Say this out loud}]
\small ``Properly de-identified'' is a \textbf{risk-based standard}, not a name-stripping checklist.
\vskip3pt
In small classes the bar is \textbf{genuinely hard to clear}. Treat it as a judgment call.
\end{tcolorbox}
\end{column}
\end{columns}
\end{frame}

%% ====== SLIDE 8: WHAT ED HAS / HASN'T SAID ======
\begin{frame}{What the Dept of Education Has --- and Hasn't --- Said}
\begin{columns}[T]
\begin{column}{0.50\textwidth}
\begin{tcolorbox}[keybox, title={\faGavel\quad No AI-specific rules}]
\small \textbf{No AI-specific FERPA regulations exist} (as of June 2026). The existing \textbf{technology-neutral} guidance governs AI vendors directly.
\end{tcolorbox}
\vskip4pt
\begin{tcolorbox}[quotebox, top=3pt, bottom=3pt]
\scriptsize The framework (school-official exception, direct control, data-minimization) is tech-neutral --- AI tools are just a subset of ``online educational services.''
\end{tcolorbox}
\end{column}
\begin{column}{0.46\textwidth}
\begin{tcolorbox}[accentbox]
\small
\textbf{The operative ED documents:}
\vskip3pt
\begin{itemize}\setlength{\itemsep}{2pt}
  \item Third-Party Service Providers under FERPA (2015)
  \item Online Educational Services: Requirements \& Best Practices (2014)
  \item \textbf{Model Terms of Service checklist (2016)} --- ready-made for vetting a tool's ToS
\end{itemize}
\end{tcolorbox}
\vskip3pt
\cmcite{studentprivacy.ed.gov}
\end{column}
\end{columns}
\end{frame}

%% ====== SECTION DIVIDER: PRACTICE ======
{
\setbeamertemplate{footline}{}
\setbeamercolor{background canvas}{bg=SlateTeal}
\begin{frame}[plain, noframenumbering]
\centering\vfill
{\fontsize{22}{28}\bfseries\selectfont\color{white} Part 2}\\[6pt]
{\color{UVMGold}\rule{5cm}{0.8pt}}\\[6pt]
{\normalsize\color{UVMGold} Practical Do's \& Don'ts}\\[4pt]
{\small\color{PaleSage} For Economics Faculty}
\vfill
\end{frame}
}

%% ====== SLIDE 9: DO / DON'T ======
\begin{frame}{Do's \& Don'ts}
\begin{columns}[T]
\begin{column}{0.48\textwidth}
\begin{tcolorbox}[dontbox, title={\faTimes\quad DON'T}]
\small
\begin{itemize}\setlength{\itemsep}{3pt}
  \item Paste \textbf{identifiable} essays, exams, or grades into a \textbf{consumer} AI account
  \item Assume \textbf{deleting the name} makes an essay safe
  \item Upload a \textbf{gradebook / roster} to ``summarize''
  \item Build a chatbot \textbf{trained on} identifiable student work without approval + a DPA
\end{itemize}
\end{tcolorbox}
\end{column}
\begin{column}{0.48\textwidth}
\begin{tcolorbox}[dobox, title={\faCheck\quad DO}]
\small
\begin{itemize}\setlength{\itemsep}{3pt}
  \item Use your institution's \textbf{vetted, contracted} AI tool for identifiable data
  \item \textbf{De-identify first} --- for re-ID risk, not just names
  \item Run a tool's ToS through ED's \textbf{2016 Model ToS checklist}
  \item Check your \textbf{institution's own AI policy} (often stricter)
  \item Ask your \textbf{registrar / privacy office} when unsure
\end{itemize}
\end{tcolorbox}
\end{column}
\end{columns}
\end{frame}

%% ====== SLIDE 10: CLASSROOM SCENARIOS ======
\begin{frame}{Concrete Classroom Scenarios}
\vskip-2pt
{\small
\begin{tabular}{>{\raggedright\arraybackslash}m{7.3cm} >{\centering\arraybackslash}m{1.3cm} >{\raggedright\arraybackslash}m{3.4cm}}
\rowcolor{SlateTeal}
\textcolor{white}{\textbf{Scenario}} & \textcolor{white}{\textbf{Verdict}} & \textcolor{white}{\textbf{Why}} \\[2pt]
\rowcolor{PaleSage}
Named essay into consumer ChatGPT for feedback & \textcolor{SoftCoral}{\faTimes} & Unconsented disclosure \\[3pt]
\rowcolor{white}
Same essay, \textbf{fully de-identified}, low re-ID risk, general tool & \textcolor{UVMGold}{\faExclamationTriangle} & Defensible --- your judgment \\[3pt]
\rowcolor{PaleSage}
Same essay, in the \textbf{university's contracted} tool & \textcolor{GoGreen}{\faCheck} & Vendor is a school official \\[3pt]
\rowcolor{white}
Upload the class gradebook to summarize performance & \textcolor{SoftCoral}{\faTimes} & Grades are PII \\[3pt]
\rowcolor{PaleSage}
Course chatbot \textbf{trained on} student submissions & \textcolor{SoftCoral}{\faTimes} & Needs approval + DPA \\[3pt]
\rowcolor{white}
AI to draft a \textbf{generic} rubric or lecture (no student data) & \textcolor{GoGreen}{\faCheck} & No education records involved \\[2pt]
\end{tabular}
}
\vskip5pt
\begin{tcolorbox}[quotebox, width=12.5cm, top=3pt, bottom=3pt]
\small The dividing line is almost always: \textbf{identifiable student data + a tool that isn't institutionally controlled.}
\end{tcolorbox}
\end{frame}

%% ====== SLIDE 11: THE MODEL-TRAINING QUESTION ======
\begin{frame}{The Question Everyone Asks: ``Does It Train on My Data?''}
\begin{columns}[T]
\begin{column}{0.50\textwidth}
\begin{tcolorbox}[keybox, title={\faBrain\quad The FERPA answer}, top=3pt, bottom=3pt]
\small
Under the school-official exception, a vendor may use PII \textbf{only for the contracted purpose} and may not reuse it for its own ends.
\vskip3pt
So training a vendor's models on \textbf{identifiable} student records is an \textbf{unauthorized reuse} --- unless you authorize it.
\vskip3pt
\textbf{De-identified} data is outside FERPA --- training on it is fine.
\end{tcolorbox}
\vskip3pt
\cmcite{ED 2014/2016 guidance; CRS R46799. A sound inference from the use-limitation rule --- ED has no AI-specific ruling.}
\end{column}
\begin{column}{0.46\textwidth}
\begin{tcolorbox}[coralbox, title={\faExclamationTriangle\quad The vendor reality}, top=2pt, bottom=2pt]
\small
\begin{itemize}\setlength{\itemsep}{1pt}
  \item OpenAI, MS 365 Copilot, and Claude (\textbf{commercial/edu} tiers) all state they don't train on your prompts \emph{by default}.
  \item These are \textbf{vendor statements}, not FERPA rulings --- and consumer tiers are excluded.
  \item Only \textbf{Microsoft} publishes an explicit FERPA ``school-official'' commitment.
\end{itemize}
\end{tcolorbox}
\vskip3pt
\cmcite{Full three-way comparison on the final (reference) slide.}
\end{column}
\end{columns}
\end{frame}

%% ====== SLIDE 12: ENFORCEMENT REALITY ======
\begin{frame}{So How Worried Should You Be? The Enforcement Reality}
\begin{columns}[T]
\begin{column}{0.50\textwidth}
\begin{tcolorbox}[keybox, title={\faBalanceScale\quad How FERPA is enforced}, top=3pt, bottom=3pt]
\small
\begin{itemize}\setlength{\itemsep}{2pt}
  \item \textbf{No private right of action} --- a student cannot sue you under FERPA \cmcite{(\textit{Gonzaga v. Doe}, 2002)}.
  \item ED's ultimate stick is \textbf{withdrawing federal funds} --- never actually used.
  \item Real ed-tech enforcement has run through \textbf{COPPA/FTC}, not FERPA (Edmodo, \$6M, 2023).
\end{itemize}
\end{tcolorbox}
\end{column}
\begin{column}{0.46\textwidth}
\begin{tcolorbox}[accentbox]
\small
\textbf{So the real exposure is \textbf{institutional and reputational}}, not a lawsuit.
\vskip3pt
That's exactly why your university \textbf{vets and contracts} tools centrally --- and why the safe move is to \textbf{use the approved one}.
\end{tcolorbox}
\vskip3pt
\begin{tcolorbox}[quotebox, top=2pt, bottom=2pt]
\scriptsize ED's only recent vendor action (the May 2026 Canvas/Instructure letter) was about a \textbf{data breach} --- not AI. There is still \textbf{no AI-specific FERPA guidance}.
\end{tcolorbox}
\end{column}
\end{columns}
\end{frame}

%% ====== SLIDE 12: TAKEAWAYS ======
\begin{frame}{Takeaways: A Faculty Checklist}
\vskip2pt
\begin{tcolorbox}[keybox, title={\faClipboardCheck\quad Before you put student data in any AI tool}]
\small
\begin{enumerate}\setlength{\itemsep}{4pt}
  \item Is it \textbf{identifiable} student data? (Remember: PII is broader than names.)
  \item If yes --- is the tool \textbf{institutionally vetted / contracted}? If not, \textbf{stop}.
  \item Or can you \textbf{properly de-identify} it (risk-based, hard in small classes)?
  \item When in doubt, run the ToS through ED's \textbf{Model ToS checklist} and ask your \textbf{privacy office}.
\end{enumerate}
\end{tcolorbox}
\vskip6pt
\begin{tcolorbox}[accentbox, width=12.5cm]
\centering\small \textbf{The one thing to remember:} approved tool \emph{or} de-identify --- never identifiable student data in a personal AI account.
\end{tcolorbox}
\end{frame}

%% ====== REFERENCE SLIDE: VENDOR COMPARISON (for Q&A / backup) ======
\begin{frame}{Reference: The Three Enterprise Tools, Compared}
\vskip-2pt
{\footnotesize
\begin{tabular}{>{\raggedright\arraybackslash}m{3.9cm} >{\centering\arraybackslash}m{2.7cm} >{\centering\arraybackslash}m{2.7cm} >{\centering\arraybackslash}m{2.7cm}}
\rowcolor{SlateTeal}
\textcolor{white}{\textbf{}} & \textcolor{white}{\textbf{ChatGPT Edu / Enterprise}} & \textcolor{white}{\textbf{MS 365 Copilot}} & \textcolor{white}{\textbf{Claude Edu / Enterprise}} \\[1pt]
\rowcolor{PaleSage}
Trains on your data by default? & \textcolor{GoGreen}{\textbf{No}} & \textcolor{GoGreen}{\textbf{No}} & \textcolor{GoGreen}{\textbf{No}} \\[1pt]
\rowcolor{white}
Consumer tier trains by default? & Yes (opt-out) & N/A & \textbf{Yes} since Aug 2025 \\[1pt]
\rowcolor{PaleSage}
Explicit FERPA ``school official'' language? & \textcolor{SoftCoral}{No} & \textcolor{GoGreen}{\textbf{Yes}} (99.33(a)) & \textcolor{SoftCoral}{No} (processor + DPA) \\[1pt]
\rowcolor{white}
Data Processing Addendum? & Yes & Yes (OST DPA) & Yes (auto-incl.) \\[1pt]
\rowcolor{PaleSage}
Retention & Configurable & Configurable & Until deleted (API: 30d) \\[1pt]
\rowcolor{white}
Zero-data-retention option? & Via API & --- & \textbf{API only, not chat} \\[1pt]
\end{tabular}
}
\vskip4pt
\begin{tcolorbox}[quotebox, width=13cm, top=2pt, bottom=2pt]
\scriptsize All entries are \textbf{vendor representations}, not adjudicated FERPA compliance --- terms change fast. \textbf{Microsoft} alone publishes an explicit FERPA school-official commitment; the others rely on the contract/DPA + institutional vetting. \cmcite{Sources: openai.com, learn.microsoft.com, privacy.claude.com --- as of June 2026.}
\end{tcolorbox}
\end{frame}

\end{document}
